\ProvidesFile{LocalityCode.tex}[Код для раздела Полиморфизм и локальность]


\newcommand{\codewidth}{5cm}

\newsavebox\funch
\begin{lrbox}{\funch}%
\begin{minipage}[\baselineskip]{\codewidth}
\texttt{function.h}
\begin{cppcode}[numbers = none]
#pragma once
#include "A.h"


void func(A&);

\end{cppcode}
\end{minipage}%
\end{lrbox}

\newsavebox\funccpp
\begin{lrbox}{\funccpp}%
\begin{minipage}[\baselineskip]{\codewidth}
\texttt{function.cpp}
\begin{cppcode}[numbers = none]
#include "function.h"

#include <iostream>

void func(A& a) {
  int x = a.f();
  int y = a.g(x);
  std::cout << y;
}
\end{cppcode}
\end{minipage}%
\end{lrbox}

\newsavebox\funcha
\begin{lrbox}{\funcha}%
\begin{minipage}[\baselineskip]{\codewidth}
\texttt{function.h}
\begin{cppcode}[numbers = none]
#pragma once
#include "???.h"


void func(???&);

\end{cppcode}
\end{minipage}%
\end{lrbox}

\newsavebox\funccppa
\begin{lrbox}{\funccppa}%
\begin{minipage}[\baselineskip]{\codewidth}
\texttt{function.cpp}
\begin{cppcode}[numbers = none]
#include "function.h"

#include <iostream>


void func(???& a) {
  int x = a.f();
  int y = a.g(x);
  std::cout << y;
}

\end{cppcode}
\end{minipage}%
\end{lrbox}

\newsavebox\funchb
\begin{lrbox}{\funchb}%
\begin{minipage}[\baselineskip]{\codewidth}
\texttt{function.h}
\begin{cppcode}[numbers = none]
#pragma once
#include "I.h"


void func(I&);

\end{cppcode}
\end{minipage}%
\end{lrbox}

\newsavebox\funccppb
\begin{lrbox}{\funccppb}%
\begin{minipage}[\baselineskip]{\codewidth}
\texttt{function.cpp}
\begin{cppcode}[numbers = none]
#include "function.h"

#include <iostream>

void func(I& a) {
  int x = a.f();
  int y = a.g(x);
  std::cout << y;
}
\end{cppcode}
\end{minipage}%
\end{lrbox}

\newsavebox\funccppc
\begin{lrbox}{\funccppc}%
\begin{minipage}[\baselineskip]{\codewidth}
\texttt{function.cpp}
\begin{cppcode}[numbers = none]
#include "function.h"

#include <iostream>


void func(I& a) {
  int x = a.f();
  int y = a.g(x);
  std::cout << y;
}

\end{cppcode}
\end{minipage}%
\end{lrbox}

\newsavebox\main
\begin{lrbox}{\main}%
\begin{minipage}[\baselineskip]{\codewidth}
\texttt{main.cpp}
\begin{cppcode}[numbers = none]
#include "function.h"

int main() {
  A a;
  func(a);
  A b;
  func(b);
  return 0;
}
\end{cppcode}
\end{minipage}%
\end{lrbox}

\newsavebox\maina
\begin{lrbox}{\maina}%
\begin{minipage}[\baselineskip]{\codewidth}
\texttt{main.cpp}
\begin{cppcode}[numbers = none]
#include "function.h"
#include "A.h"
#include "B.h"

int main() {
  A a;
  func(a);
  B b;
  func(b);
  return 0;
}
\end{cppcode}
\end{minipage}%
\end{lrbox}

\newsavebox\mainb
\begin{lrbox}{\mainb}%
\begin{minipage}[\baselineskip]{\codewidth}
\texttt{main.cpp}
\begin{cppcode}[numbers = none]
#include "function.h"
#include "A.h"
#include "BI.h"

int main() {
  A a;
  func(a);
  B b;
  BI bi = BI{&b};
  func(bi);
}
\end{cppcode}
\end{minipage}%
\end{lrbox}

\newsavebox\mainc
\begin{lrbox}{\mainc}%
\begin{minipage}[\baselineskip]{\codewidth}
\texttt{main.cpp}
\begin{cppcode}[numbers = none]
#include "function.h"

int main() {
  A a;
  func(a);
  A b;
  func(b);
  return 0;
}
\end{cppcode}
\end{minipage}%
\end{lrbox}

\newsavebox\classAh
\begin{lrbox}{\classAh}%
\begin{minipage}[\baselineskip]{\codewidth}
\texttt{A.h}
\begin{cppcode}[numbers = none]
#pragma once

struct A {
  int f() const;
  int g(int) const;
};
\end{cppcode}
\end{minipage}%
\end{lrbox}

\newsavebox\classAcpp
\begin{lrbox}{\classAcpp}%
\begin{minipage}[\baselineskip]{\codewidth}
\texttt{A.cpp}
\begin{cppcode}[numbers = none]
#include "A.h"

int A::f() const {
  return /*something*/;
}

int A::g(int x) const {
  return /*something of x*/;
}
\end{cppcode}
\end{minipage}%
\end{lrbox}

\newsavebox\classBh
\begin{lrbox}{\classBh}%
\begin{minipage}[\baselineskip]{\codewidth}
\texttt{B.h}
\begin{cppcode}[numbers = none]
#pragma once

struct B {
  int f() const;
  int g(int) const;
};
\end{cppcode}
\end{minipage}%
\end{lrbox}

\newsavebox\classBcpp
\begin{lrbox}{\classBcpp}%
\begin{minipage}[\baselineskip]{\codewidth}
\texttt{B.cpp}
\begin{cppcode}[numbers = none]
#include "B.h"

int B::f() const {
  return /*something*/;
}

int B::g(int x) const {
  return /*something of x*/;
}
\end{cppcode}
\end{minipage}%
\end{lrbox}

\newsavebox\classIh
\begin{lrbox}{\classIh}%
\begin{minipage}[\baselineskip]{\codewidth}
\texttt{I.h}
\begin{cppcode}[numbers = none]
#pragma once

struct I {
  int f() const;
  int g(int) const;
};
\end{cppcode}
\end{minipage}%
\end{lrbox}

\newsavebox\classAIh
\begin{lrbox}{\classAIh}%
\begin{minipage}[\baselineskip]{\codewidth}
\texttt{A.h}
\begin{cppcode}[numbers = none]
#pragma once

struct A : I {
  int f() const override;
  int g(int) const override;
};
\end{cppcode}
\end{minipage}%
\end{lrbox}

\newsavebox\classAIcpp
\begin{lrbox}{\classAIcpp}%
\begin{minipage}[\baselineskip]{\codewidth}
\texttt{A.cpp}
\begin{cppcode}[numbers = none]
#include "A.h"

int A::f() const {
  return /*something*/;
}

int A::g(int x) const {
  return /*something of x*/;
}
\end{cppcode}
\end{minipage}%
\end{lrbox}

\newsavebox\classBIh
\begin{lrbox}{\classBIh}%
\begin{minipage}[\baselineskip]{\codewidth}
\texttt{B.h}
\begin{cppcode}[numbers = none]
#pragma once

struct B : I {
  int f() const override;
  int g(int) const override;
};
\end{cppcode}
\end{minipage}%
\end{lrbox}

\newsavebox\classBIcpp
\begin{lrbox}{\classBIcpp}%
\begin{minipage}[\baselineskip]{\codewidth}
\texttt{B.cpp}
\begin{cppcode}[numbers = none]
#include "B.h"

int B::f() const {
  return /*something*/;
}

int B::g(int x) const {
  return /*something of x*/;
}
\end{cppcode}
\end{minipage}%
\end{lrbox}

\newsavebox\classIVh
\begin{lrbox}{\classIVh}%
\begin{minipage}[\baselineskip]{\codewidth+0.5cm}
\texttt{I.h}
\begin{cppcode}[numbers = none]
#pragma once

struct I {
  virtual int f() const = 0;
  virtual int g(int) const = 0;
};
\end{cppcode}
\end{minipage}%
\end{lrbox}

\newsavebox\classBIAh
\begin{lrbox}{\classBIAh}%
\begin{minipage}[\baselineskip]{\codewidth}
\texttt{BI.h}
\begin{cppcode}[numbers = none]
#pragma once
#include "I.h"
#include "B.h"

struct BI : I {
  int f() const override;
  int g(int) const override;
  B* b;
};
\end{cppcode}
\end{minipage}%
\end{lrbox}

\newsavebox\classBIAcpp
\begin{lrbox}{\classBIAcpp}%
\begin{minipage}[\baselineskip]{\codewidth}
\texttt{BI.cpp}
\begin{cppcode}[numbers = none]
#include "BI.h"

int B::f() const {
  return b->f();
}

int B::g(int x) const {
  return b->g(x);
}
\end{cppcode}
\end{minipage}%
\end{lrbox}
